\input{preamble}

\begin{document}

\title{Linear Algebra}
\maketitle

\phantomsection
\label{section-phantom}

\tableofcontents

\section{Signature}
\label{section-signature}

The signature of a form quadratic form can be
equivalently computed as the
maximal dimension of an isotropic subspace,
or as the maximal dimension of the
subspace where the form is negative.

\begin{remark}
\label{remark-nondegenerate-vs-anisotropic}
Let $B:V \times V \to k$ be a bilinear form.
The {\it radical} of $B$ is
$$
\operatorname{rad}(B)=
\{v \in V : B(v,w)=0\text{ for all }w
\in V\}.
$$
The form is {\it nondegenerate} if
$\operatorname{rad}(B)=0$. In coordinates,
if $A$ is the Gram matrix of $B$, this is
equivalent to $\det(A)\neq 0$.

This should not be confused with the
condition that there are no nonzero isotropic
vectors. A vector $v\neq 0$ is {\it
isotropic} if $B(v,v)=0$, and a form with no
nonzero isotropic vectors is called
{\it anisotropic}. Anisotropic implies
nondegenerate, since a nonzero vector in the
radical would be isotropic. The converse is
false; for example the form with matrix
$$
\begin{pmatrix}
1&0\\
0&-1
\end{pmatrix}
$$
is nondegenerate, but $(1,1)$ is isotropic.
\end{remark}

\section{Adjoint}
\label{section-adjoint}

\begin{definition}
\label{definition-adjoint}
Let  $A:(V,\left\langle{}\cdot,\cdot\right\rangle{}_V) \to
(W,\left\langle{}\cdot,\cdot\right\rangle{}_W)$ be a
linear transformation of inner product vector
spaces, its {\it adjoint} is the
map
$$
A^{\dag}:W\to V
$$
satisfying
\begin{equation}
\label{equation-adjoint}
\left\langle{}Av,w\right\rangle{}_W=\left\langle{}v,A^\dag w\right\rangle{}_V
\end{equation}
\end{definition}

Let $G$ be a matrix representation for
$\left\langle{}\cdot,\cdot\right\rangle{}_V$ and $H$ for
$\left\langle{}\cdot,\cdot\right\rangle{}_W$. Then Eq.
\ref{equation-adjoint} is
$$
v^{\mathbf{T}}A^{\mathbf{T}}Hw=(Av)
^{\mathbf{T}}Hw=v^{\mathbf{T}}GA^\dag
w
$$
So
\begin{equation}
\label{equation-adjoint-transpose}
A^{\mathbf{T}}H=GA^\dag \iff
A^\dag:=G^{-1}A^{\mathbf{T}}H
\end{equation}
which gives the relationship between the
usual transpose matrix and the adjoint.

\begin{definition}
\label{definition-Frobenius-norm}
In the space of matrices we define {\it
Frobenius norm} as
\begin{equation}
\label{equation-Frobenius-norm}
\|A\|^2=\text{tr}A^\dag
A=\text{tr}(G^{-1}A^{\mathbf{T}}HA
\end{equation}
\end{definition}

\bibliography{my}
\bibliographystyle{amsalpha}

\end{document}
